\documentclass[a4paper,11pt]{article}

% --- Packages ---
\usepackage[utf8]{inputenc}
\usepackage{geometry}
\geometry{left=2.5cm, right=2.5cm, top=2.5cm, bottom=2.5cm}
\usepackage{xcolor}
\usepackage{hyperref}
\usepackage{listings}
\usepackage{graphicx}
\usepackage{titlesec}
\usepackage{parskip}
\usepackage{amsmath}

% --- Colors ---
\definecolor{codegray}{rgb}{0.95,0.95,0.95}
\definecolor{codeblue}{rgb}{0.1,0.1,0.5}
\definecolor{ethblue}{RGB}{31, 64, 122}

% --- Hyperlinks ---
\hypersetup{
    colorlinks=true,
    linkcolor=ethblue,
    filecolor=magenta,      
    urlcolor=cyan,
}

% --- Code Style ---
\lstset{
    backgroundcolor=\color{codegray},
    basicstyle=\ttfamily\small,
    breakatwhitespace=false,         
    breaklines=true,                 
    captionpos=b,                    
    keepspaces=true,                 
    showspaces=false,                
    showstringspaces=false,
    showtabs=false,                  
    tabsize=2,
    frame=single,
    rulecolor=\color{black!30}
}

% --- Title Info ---
\title{\textbf{\Huge ExciView} \\ \Large BSE Wavefunction Analysis Toolkit for FHI-aims}
\author{\textbf{Ruiyi Zhou} \\ ETH Zurich \\ \texttt{ruiyi.zhou@phys.chem.ethz.ch}}
\date{\today}

\begin{document}

\maketitle
\tableofcontents
\newpage

% =============================================================================
\section{Introduction}
% =============================================================================
\textbf{ExciView} is a post-processing toolkit designed to analyze many-body excitonic wavefunctions calculated using the Bethe-Salpeter Equation (BSE) formalism in the FHI-aims software package. 

Because BSE Hamiltonians are computationally expensive, ExciView employs an optimized I/O strategy to bridge the gap between:
\begin{itemize}
    \item \textbf{Reciprocal Space:} Eigenvector weights ($W_{\mathbf{k}}$) and band transitions.
    \item \textbf{Real Space:} Atomic orbital projections (Mulliken) and volumetric densities (Cube files).
\end{itemize}

% =============================================================================
\section{Installation}
% =============================================================================
ExciView is a Python package that requires \texttt{numpy}, \texttt{scipy}, and \texttt{ase} (Atomic Simulation Environment).

\subsection{Prerequisites}
Ensure you have Python 3.8+ installed. 
\begin{lstlisting}[language=bash]
pip install numpy scipy ase
\end{lstlisting}

\subsection{Installing ExciView}
Navigate to the root directory of the package (where \texttt{pyproject.toml} is located) and run:
\begin{lstlisting}[language=bash]
pip install .
\end{lstlisting}

Once installed, you can launch the toolkit from any directory containing your calculation data by typing:
\begin{lstlisting}[language=bash]
exciview
\end{lstlisting}

% =============================================================================
\section{The Main Menu (Workflow Overview)}
% =============================================================================
When you launch ExciView, you will see the interactive menu. The workflow is divided into four logical pipelines.

\begin{lstlisting}
==================================================
          BSE ANALYSIS TOOLKIT (ExciView)         
==================================================
1. Phase 1: Generate Mulliken Inputs
2. Phase 2: Analyze Mulliken Output
3. Phase 3: Generate Cube Inputs (Avg Density)
4. Phase 4: Analyze Cube Files (Avg Density)
5. Phase 5: BZ & Band Analysis (Text Only)
6. Phase 6: Generate Conditional Inputs
7. Phase 7: Analyze Conditional Density
0. Exit
\end{lstlisting}

\subsection{Workflow Logic}
Most analysis tasks follow a \textbf{two-step process} to save computational resources:
\begin{enumerate}
    \item \textbf{Generation Phase (Odd Numbers):} The script reads the BSE binary, identifies the most important k-points/bands (using a threshold), and generates a \texttt{snippet.in} file. You copy this into \texttt{control.in} and run FHI-aims.
    \item \textbf{Analysis Phase (Even Numbers):} The script reads the output files generated by FHI-aims and combines them to produce the final result.
\end{enumerate}

\newpage

% =============================================================================
\section{Detailed Usage Guide}
% =============================================================================

\subsection{Input Requirements}
Before starting, ensure you have the following files in your directory:
\begin{itemize}
    \item \texttt{BSE\_eigenvectors.dat}: The binary output from ELSI/FHI-aims.
    \item \texttt{geometry.in}: (Optional) Used to determine lattice vectors for physical unit conversion.
\end{itemize}
You will need to know:
\begin{itemize}
    \item \textbf{State Index:} Which exciton state to analyze (0-based index).
    \item \textbf{Nv / Nc:} Number of Valence and Conduction bands used in the BSE calculation.
    \item \textbf{Grid:} The k-grid dimensions ($N_{kx}, N_{ky}, N_{kz}$).
\end{itemize}

% -----------------------------------------------------------------------------
\subsection{Pipeline A: Atomic Population (Mulliken)}
\textit{Use this to determine which atoms the hole and electron reside on.}

\textbf{Option 1: Generate Inputs} \\
Generates \texttt{mulliken\_snippet.in}. Copy this content to \texttt{control.in}.

\textbf{Option 2: Analyze Output} \\
Reads the \texttt{bandmlk*.out} files.
\begin{itemize}
    \item \textbf{Pattern:} The naming format of output files. Default: \texttt{bandmlk\{\}.out}.
    \item \textbf{Offset:} The number of the \textit{first} file generated. FHI-aims usually starts numbering outputs at \texttt{1001} or \texttt{1002}. Check your folder to see the first number.
\end{itemize}
\textbf{Output:} \texttt{exciton\_analysis\_state\_X.dat} (Contains s/p/d/f breakdown per atom).

% -----------------------------------------------------------------------------
\subsection{Pipeline B: Average Volumetric Density}
\textit{Use this to visualize the general shape of the hole and electron clouds.}

\textbf{Option 3: Generate Inputs} \\
Generates \texttt{cube\_snippet.in}. It requests \texttt{eigenstate\_density} ($|\psi|^2$) only for bands with significant weight.

\textbf{Option 4: Sum Cubes} \\
Reads the \texttt{.cube} files and performs a weighted incoherent sum:
\[ \rho_{avg}(\mathbf{r}) = \sum_{k,n} W_{k,n} |\psi_{n,k}(\mathbf{r})|^2 \]
\textbf{Output:} \texttt{avg\_hole\_state\_X.cube} and \texttt{avg\_elec\_state\_X.cube}.

% -----------------------------------------------------------------------------
\subsection{Pipeline C: Reciprocal Space Analysis}
\textit{Use this for a quick statistical overview without running FHI-aims again.}

\textbf{Option 5: BZ \& Band Analysis} \\
Prints a report to the terminal showing:
\begin{itemize}
    \item Dominant k-points (coordinates and weights).
    \item Dominant orbital transitions (e.g., HOMO $\to$ LUMO+1).
    \item Estimated exciton radius (based on k-space spread).
\end{itemize}

% -----------------------------------------------------------------------------
\subsection{Pipeline D: Conditional Density}
\textit{Use this to visualize the electron distribution for a fixed hole position.}

\textbf{Option 6: Generate Inputs} \\
Generates \texttt{cube\_cond\_snippet.in}. Requests complex wavefunctions (\texttt{eigenstate} and \texttt{eigenstate\_imag}).

\textbf{Option 7: Analyze Density} \\
Requires a fixed hole coordinate $\mathbf{r}_h$. Performs a coherent summation:
\[ \Psi_{cond}(\mathbf{r}_e) = \sum_{c, \mathbf{k}} \left( \sum_v A_{vc}^{\mathbf{k}} \psi_{v\mathbf{k}}^*(\mathbf{r}_h) \right) \psi_{c\mathbf{k}}(\mathbf{r}_e) \]
\textbf{Output:} \texttt{cond\_density\_state\_X.cube}.

\newpage

% =============================================================================
\section{Troubleshooting}
% =============================================================================

\textbf{Q: What is the "Offset" in Phase 2/4?} \\
A: FHI-aims numbers output files sequentially based on the order of commands in \texttt{control.in}. If your first \texttt{output band\_mulliken} command results in a file named \texttt{bandmlk1011.out}, then your offset is \textbf{1011}.

\textbf{Q: My Conditional Density (Phase 7) looks wrong/empty.} \\
A: Ensure your FHI-aims version supports \texttt{output cube eigenstate\_imag}. If it does not, you cannot retrieve the phase information necessary for periodic systems ($k \neq \Gamma$).

\textbf{Q: "File Not Found" errors.} \\
A: Ensure you have run FHI-aims after the Generation Phase (Odd numbers) and before the Analysis Phase (Even numbers). The Python script does not run FHI-aims for you; it only prepares the input and analyzes the output.

\end{document}